% GENERATED FILE. Edit canonical lessons, then run npm run generate:books.
% canonical-source-hash: fnv1a64:8c8be2c4274a1154
% canonical-lessons: RU-C31-back, RU-C31-tongue, RU-C31-lip, RU-C31-face, RU-C31-blood

\chapter{Back, Tongue, Lip, Face, Blood}
\label{ch:ru-back-tongue-lip-face-blood}
\hlchaptermodality{31}

\begin{chapteropening}
\textbf{By the end of this chapter:} \emph{I can name the back, the tongue, a lip, the face and blood.}

Complete the last lesson of chapter 31: I can name the back, the tongue, a lip, the face and blood.
\end{chapteropening}

\section[\texorpdfstring{\ru{спина} (spiná)}{спина (spiná)}]{\ru{спина} (spiná) — the back}

\label{lesson:RU-C31-back}



\begin{quote}
Before the new one: say the Russian for a shoulder, then the Russian for a knee.
\end{quote}

\subsection*{You'll want to know: \ru{спина}}
\textbf{\ru{спина}} (\emph{spiná}) — \textquotedblleft{}the back\textquotedblright{}.

\begin{cousinweb}
\textbf{\ru{спина}} — origin disputed. One view takes it from Latin \emph{spīna}, \textquotedblleft{}thorn, backbone\textquotedblright{}, through Old Polish \emph{spina}, \textquotedblleft{}backbone\textquotedblright{}; Vasmer also allows a native word, a cousin of the Latin one and of Latvian \emph{spina}, \textquotedblleft{}a switch\textquotedblright{}. It is first recorded in the 1600s (Vasmer).
\end{cousinweb}

\begin{grammarlens}[title={What you've built}]
One more word for this chapter.
\end{grammarlens}

\subsection*{Your turn}
\begin{itemize}
\raggedright
  \item \emph{Say it:} \emph{spiná}
  \item \emph{Say it:} \emph{spiná}, once more
  \item \emph{Say it:} say \emph{spiná}
  \item \emph{Recall:} say the Russian for a shoulder, then the Russian for a knee, then say \emph{spiná} again
\end{itemize}

\begin{tcolorbox}[breakable,colback=teal!4,colframe=teal!35!black,title={Before you move on}]
What does \ru{спина} mean? (\textquotedblleft{}The back\textquotedblright{}.) Say it once more.
\end{tcolorbox}
\section[\texorpdfstring{\ru{язык} (yazýk)}{язык (yazýk)}]{\ru{язык} (yazýk) — the tongue; a language}

\label{lesson:RU-C31-tongue}



\begin{quote}
Before the new one: say the Russian for a knee, then the Russian for the back.
\end{quote}

\subsection*{You'll want to know: \ru{язык}}
\textbf{\ru{язык}} (\emph{yazýk}) — \textquotedblleft{}the tongue; a language\textquotedblright{}. Both meanings, as in English \textquotedblleft{}tongue\textquotedblright{}.

\begin{cousinweb}
\textbf{\ru{язык}} $\leftarrow$ Proto-Slavic \emph{*językŭ}, \textquotedblleft{}tongue\textquotedblright{} — inherited, from the old Indo-European word for \textquotedblleft{}tongue\textquotedblright{} that also gave Old Latin \emph{dingua}, later \emph{lingua}, and English \emph{tongue} (Vasmer). The Slavic shape is irregular. \textquotedblleft{}Language\textquotedblright{} grew from the organ, as in English \emph{tongue}.
\end{cousinweb}

\begin{grammarlens}[title={What you've built}]
One more word for this chapter.
\end{grammarlens}

\subsection*{Your turn}
\begin{itemize}
\raggedright
  \item \emph{Say it:} \emph{yazýk}
  \item \emph{Say it:} \emph{yazýk}, once more
  \item \emph{Say it:} say \emph{yazýk}
  \item \emph{Recall:} say the Russian for a knee, then the Russian for the back, then say \emph{yazýk} again
\end{itemize}

\begin{tcolorbox}[breakable,colback=teal!4,colframe=teal!35!black,title={Before you move on}]
What does \ru{язык} mean? (\textquotedblleft{}The tongue; a language\textquotedblright{}.) Say it once more.
\end{tcolorbox}
\section[\texorpdfstring{\ru{губа} (gubá)}{губа (gubá)}]{\ru{губа} (gubá) — a lip}

\label{lesson:RU-C31-lip}



\begin{quote}
Before the new one: say the Russian for the back, then the Russian for the tongue; a language.
\end{quote}

\subsection*{You'll want to know: \ru{губа}}
\textbf{\ru{губа}} (\emph{gubá}) — \textquotedblleft{}a lip\textquotedblright{}.

\begin{cousinweb}
\textbf{\ru{губа}} $\leftarrow$ Proto-Slavic \emph{*gǫba} — inherited; Ukrainian \emph{\ru{губа}} and Polish \emph{gęba} mean \textquotedblleft{}mouth\textquotedblright{} (Vasmer). The deeper origin is unclear: Vasmer weighs Greek \emph{gamphai}, \textquotedblleft{}jaws\textquotedblright{}, and a link with the \emph{\ru{губа}} that means \textquotedblleft{}sponge\textquotedblright{}, without deciding.
\end{cousinweb}

\begin{grammarlens}[title={What you've built}]
One more word for this chapter.
\end{grammarlens}

\subsection*{Your turn}
\begin{itemize}
\raggedright
  \item \emph{Say it:} \emph{gubá}
  \item \emph{Say it:} \emph{gubá}, once more
  \item \emph{Say it:} say \emph{gubá}
  \item \emph{Recall:} say the Russian for the back, then the Russian for the tongue; a language, then say \emph{gubá} again
\end{itemize}

\begin{tcolorbox}[breakable,colback=teal!4,colframe=teal!35!black,title={Before you move on}]
What does \ru{губа} mean? (\textquotedblleft{}A lip\textquotedblright{}.) Say it once more.
\end{tcolorbox}
\section[\texorpdfstring{\ru{лицо} (litsó)}{лицо (litsó)}]{\ru{лицо} (litsó) — the face}

\label{lesson:RU-C31-face}



\begin{quote}
Before the new one: say the Russian for the tongue; a language, then the Russian for a lip.
\end{quote}

\subsection*{You'll want to know: \ru{лицо}}
\textbf{\ru{лицо}} (\emph{litsó}) — \textquotedblleft{}the face\textquotedblright{}.

\begin{cousinweb}
\textbf{\ru{лицо}} $\leftarrow$ Proto-Slavic \emph{*lice}, built on \emph{*likŭ}, \textquotedblleft{}face\textquotedblright{}, the word Russian keeps as \emph{\ru{лик}}, \textquotedblleft{}a face on an icon\textquotedblright{} (Vasmer). Vasmer compares Old Irish \emph{lecco}, \textquotedblleft{}cheek\textquotedblright{}.
\end{cousinweb}

\begin{grammarlens}[title={What you've built}]
One more word for this chapter.
\end{grammarlens}

\subsection*{Your turn}
\begin{itemize}
\raggedright
  \item \emph{Say it:} \emph{litsó}
  \item \emph{Say it:} \emph{litsó}, once more
  \item \emph{Say it:} say \emph{litsó}
  \item \emph{Recall:} say the Russian for the tongue; a language, then the Russian for a lip, then say \emph{litsó} again
\end{itemize}

\begin{tcolorbox}[breakable,colback=teal!4,colframe=teal!35!black,title={Before you move on}]
What does \ru{лицо} mean? (\textquotedblleft{}The face\textquotedblright{}.) Say it once more.
\end{tcolorbox}
\section[\texorpdfstring{\ru{кровь} (krov)}{кровь (krov)}]{\ru{кровь} (krov) — blood}

\label{lesson:RU-C31-blood}



\begin{quote}
Before the new one: say the Russian for a lip, then the Russian for the face.
\end{quote}

\subsection*{You'll want to know: \ru{кровь}}
\textbf{\ru{кровь}} (\emph{krov}) — \textquotedblleft{}blood\textquotedblright{}.

\begin{cousinweb}
\textbf{\ru{кровь}} $\leftarrow$ Proto-Slavic \emph{*kry}, with the stem \emph{*krŭv-} in its other forms — inherited. Vasmer lists Lithuanian \emph{kraũjas}, \textquotedblleft{}blood\textquotedblright{}, Latin \emph{cruor}, \textquotedblleft{}gore\textquotedblright{}, Greek \emph{kréas}, \textquotedblleft{}meat\textquotedblright{}, and Old Norse \emph{hrár}, \textquotedblleft{}raw\textquotedblright{}, so English \emph{raw} is a cousin, from PIE \emph{*kréwh\textsubscript{2}-}, \textquotedblleft{}raw flesh\textquotedblright{}.
\end{cousinweb}

\begin{grammarlens}[title={What you've built}]
One more word for this chapter.
\end{grammarlens}

\subsection*{Your turn}
\begin{itemize}
\raggedright
  \item \emph{Say it:} \emph{krov}
  \item \emph{Say it:} \emph{krov}, once more
  \item \emph{Say it:} say \emph{krov}
  \item \emph{Recall:} say the Russian for a lip, then the Russian for the face, then say \emph{krov} again
\end{itemize}

\begin{tcolorbox}[breakable,colback=teal!4,colframe=teal!35!black,title={Before you move on}]
What does \ru{кровь} mean? (\textquotedblleft{}Blood\textquotedblright{}.) Say it once more.
\end{tcolorbox}
